\section{Introdução}

O câncer de mama permanece como um dos principais desafios da saúde pública contemporânea, tanto pela sua elevada incidência quanto pelo impacto social, clínico e econômico associado ao diagnóstico e ao tratamento da doença. Nesse cenário, a detecção precoce continua sendo um dos fatores mais importantes para ampliar as chances de intervenção em estágios iniciais e melhorar o prognóstico das pacientes. Entre os métodos de rastreamento disponíveis, a mamografia consolidou-se como uma das principais ferramentas para identificação de alterações suspeitas, sendo amplamente utilizada em programas de rastreamento populacional e na prática clínica. Apesar de sua relevância, a interpretação de mamografias permanece complexa, pois depende da análise de achados muitas vezes sutis, como massas, microcalcificações, distorções arquiteturais e assimetrias, além de estar sujeita à variabilidade interobservador e à alta carga de leitura dos especialistas (Abdelhafiz et al., 2019; Elías-Cabot et al., 2024).

Durante muitos anos, sistemas tradicionais de \textit{computer-aided detection} (CAD) foram desenvolvidos com o objetivo de auxiliar radiologistas na leitura mamográfica. Entretanto, a literatura mostra que tais sistemas apresentaram ganhos limitados na prática, sobretudo em razão do aumento de falsos positivos e da dificuldade de melhorar substancialmente a acurácia diagnóstica. A partir do avanço do aprendizado profundo, especialmente das Redes Neurais Convolucionais (CNNs), esse panorama começou a mudar. Segundo Abdelhafiz et al. (2019), as CNNs passaram a ser empregadas em diversas tarefas relacionadas à mamografia, incluindo detecção de lesões, localização, classificação, recuperação de imagens e avaliação de risco, destacando-se como alternativas mais robustas do que os métodos CAD convencionais.

Esse progresso também está relacionado à evolução das arquiteturas profundas no campo mais amplo da visão computacional. O trabalho de Szegedy et al. (2014), ao apresentar a arquitetura Inception/GoogLeNet, mostrou que era possível ampliar profundidade e largura de redes neurais mantendo eficiência computacional, estabelecendo um marco importante para tarefas complexas de classificação visual. Embora esse estudo não tenha sido desenvolvido para mamografia, ele representa um fundamento arquitetural importante para a consolidação de redes profundas mais eficientes, posteriormente adaptadas para aplicações em imagens médicas.

No domínio específico da mama, os estudos mais recentes indicam que os modelos de aprendizado profundo já ultrapassaram a fase de simples prova de conceito e passaram a demonstrar utilidade em cenários clinicamente relevantes. Wu et al. (2020), por exemplo, desenvolveram uma rede profunda treinada e avaliada em mais de 200 mil exames de rastreamento, totalizando mais de 1 milhão de imagens, alcançando AUC de 0,895 para predição da presença de câncer na mama. Além disso, os autores mostraram que a combinação entre a predição da rede e a leitura humana pode superar o desempenho isolado de cada uma das partes. Em complemento, Elías-Cabot et al. (2024) avaliaram o uso de inteligência artificial como suporte à dupla leitura humana em um programa populacional real com mamografia digital e tomossíntese, observando aumento na taxa de detecção de câncer e no valor preditivo positivo dos recalls. Em conjunto, esses achados reforçam a relevância da IA como ferramenta de apoio à decisão médica em rastreamento mamográfico.

Além do desempenho global dos modelos, a literatura recente também tem aprofundado o estudo de tarefas específicas e estratégias metodológicas mais refinadas. Pesapane et al. (2023) desenvolveram modelos de \textit{deep learning} voltados à detecção e classificação de microcalcificações, obtendo resultados elevados em sensibilidade, especificidade e AUC, o que evidencia a capacidade das redes neurais em lidar com achados pequenos e de alta importância clínica. Já Camurdan et al. (2025) propuseram uma abordagem baseada em \textit{patches}, \textit{curriculum learning} e anotações mistas, mostrando que é possível melhorar desempenho e explicabilidade mesmo com quantidade limitada de rótulos fortes. Em outra direção, Yoshitsugu, Kishimoto e Takemura (2026) exploraram a estratificação por presença de região de interesse (ROI), indicando que a organização do problema de classificação também influencia diretamente a qualidade dos resultados em tarefas normal/anormal e benigno/maligno.

Outro fator decisivo para o avanço da área é a disponibilidade de bases de dados amplas, heterogêneas e bem anotadas. Uma limitação recorrente da literatura em mamografia está justamente na fragmentação dos conjuntos de dados, que variam quanto ao formato, qualidade, densidade de anotações, equilíbrio entre classes e representatividade populacional. Nesse contexto, Bhole, Suba e Parekh (2025) propuseram o Mammo-Bench, um \textit{benchmark} de grande escala construído a partir da integração de sete bases distintas, reunindo 74.436 imagens de 26.500 pacientes de sete países. Os autores mostram que uma base mais ampla e padronizada contribui para o treinamento de modelos mais robustos e favorece comparações experimentais mais consistentes. Para pesquisas comparativas, esse tipo de iniciativa é particularmente importante, pois reduz a influência de vieses associados a bases pequenas, muito específicas ou fortemente desbalanceadas.

Paralelamente à consolidação das abordagens clássicas, a computação quântica passou a ser investigada como uma alternativa emergente para problemas de aprendizado de máquina. No contexto do \textit{Quantum Machine Learning} (QML), diferentes autores têm proposto modelos híbridos que combinam etapas clássicas de extração de características com circuitos quânticos parametrizados. Schetakis et al. (2022) discutem arquiteturas híbridas para classificação em conjuntos de dados ruidosos, combinando redes híbridas, circuitos paramétricos e \textit{data re-uploading}. Watkins, Chen e Yoo (2023), por sua vez, acrescentam a esse debate a questão da privacidade, propondo um modelo híbrido quântico-clássico treinado com privacidade diferencial. Tais estudos mostram que o campo quântico vem avançando conceitualmente, ainda que em cenários controlados e majoritariamente simulados.

Mais recentemente, arquiteturas híbridas voltadas a imagens passaram a explorar a integração entre operações convolucionais clássicas e componentes quânticos treináveis. Long et al. (2025), com o modelo QCQ-CNN, propõem uma arquitetura composta por um filtro quântico, um bloco clássico raso e um classificador quântico variacional, defendendo que a inclusão de parâmetros quânticos treináveis pode ampliar a expressividade do modelo. Em outra linha, Daka e Bhattacharyya (2026) apresentam o No-QCNN, uma arquitetura que procura operar com um fluxo quântico-convolucional mais direto, especialmente em cenários de poucos dados. Ainda assim, os próprios autores reconhecem que os resultados permanecem limitados por fatores como escala reduzida, sobrecarga de codificação de dados, ruído de amostragem e restrições inerentes aos dispositivos NISQ. Além disso, trabalhos como o de Egginger, Sakhnenko e Lorenz (2024) mostram que métodos quânticos, como kernels quânticos, são fortemente dependentes da escolha de hiperparâmetros e da adequação entre o \textit{feature map} e o conjunto de dados, o que reforça a necessidade de análises empíricas cuidadosas antes de qualquer alegação de vantagem prática.

Dessa forma, observa-se na literatura uma assimetria importante: enquanto os modelos clássicos baseados em CNNs já apresentam resultados maduros em mamografia, inclusive com evidências de aplicação clínica real, os modelos híbridos quântico-clássicos ainda se encontram em estágio mais exploratório, com validações predominantemente realizadas em bases pequenas, tarefas simplificadas ou ambientes simulados. Essa diferença não diminui o interesse científico das abordagens quânticas; ao contrário, evidencia a existência de uma lacuna relevante para pesquisa. Ainda faltam estudos comparativos conduzidos de maneira controlada, utilizando o mesmo problema, as mesmas classes diagnósticas e critérios experimentais equivalentes, de modo a permitir uma avaliação mais justa entre arquiteturas clássicas e híbridas.

É nesse contexto que se insere o presente trabalho. Esta pesquisa propõe uma análise comparativa entre Redes Neurais Convolucionais clássicas e arquiteturas híbridas quântico-clássicas aplicadas à classificação de mamografias digitais em três categorias diagnósticas: \textit{Normal}, \textit{Benigno} e \textit{Maligno}. O objetivo é investigar, a partir de um \textit{baseline} clássico robusto, a viabilidade técnica, o desempenho preditivo e as limitações computacionais de um modelo híbrido quântico-clássico, mantendo controle sobre as principais variáveis experimentais. Assim, o estudo busca contribuir tanto para a literatura de IA aplicada à mamografia quanto para a discussão, ainda em aberto, sobre o real potencial de integração entre computação quântica e diagnóstico por imagem.